\chapter{范畴 \texorpdfstring{$\Shv_{\mathrm{Nilp}}(\Bun_G)$}{ShvNilp(BunG)}}\label{ch:11}
\srcpage{57}
\chref{10}表述了 de Rham 范畴猜想\cite{AG15,GR24a}
\[\DMod(\Bun_G)\simeq\IndCoh_{\mathrm{Nilp}}(\LocSys^{\mathrm{dR}}_{\checkG}).\]
本章讨论其他层理论中的对应版本，特别是\cite{BZN16,NY19}中的 Betti 表述。在 Betti 情形，我们使用复解析叠 $\Bun_G(\mathbb C)$ 上的普通层。正确的自守范畴并非所有层，而是具有幂零微局部奇异支撑的满子范畴。

\section{代数叠上的所有层}\label{sec:11.1}\index{层!代数叠上的所有层|hyperpage}\index{光滑下降|hyperpage}
对 $\mathbb C$ 上有限型仿射概形 $S$，置
\[\Shv^{\mathrm{all}}(S):=\Shv^{\mathrm{all}}(S(\mathbb C)),\]
不加任何可构造性条件。若 $Y$ 为代数叠，则通过光滑下降定义 $Y$ 上的层：
\begin{definition*}{}
\[\Shv^{\mathrm{all}}(Y):=\lim_{(S\to Y)\in(\operatorname{Aff}^{\mathrm{sm}}_{/Y})^{\mathrm{op}}}\Shv^{\mathrm{all}}(S),\]
其中过渡函子为拉回。
\end{definition*}
等价地，利用伴随关系，也可以将它写成余极限
\[\Shv^{\mathrm{all}}(Y)\simeq\underset{\substack{S\to Y\\\text{光滑}}}{\operatorname{colim}}\Shv^{\mathrm{all}}(S),\]
其中过渡函子为相应的左伴随。一般的范畴原理如下：若 $L_{ij}:\mathcal C_i\rightleftarrows\mathcal C_j:R_{ij}$ 是可呈示范畴之间的一对伴随函子，则

\begin{transnote}{含糊点拨}
这条原理要求范畴可呈示：有紧生成元、滤余极限够多，否则按左伴随取的余极限与按右伴随取的极限未必对得上。它把「光滑下降给出的极限」与「按左伴随写出的余极限」等同起来，正是下面把 $\Shv^{\mathrm{all}}(Y)$ 写成余极限的依据；没有它，两种写法只是形式记号。[依据：原书 p.57 本段；并见 nLab "locally presentable category"（该条目无 Last revised，2026-11-17 访问）。]
\end{transnote}
\[\underset{L_{ij}}{\operatorname{colim}}\mathcal C_i\simeq\lim_{R_{ij}}\mathcal C_i.\]
因此，$\Shv^{\mathrm{all}}(\Bun_G)$ 是一个由 $\Bun_G$ 的解析图册构成的大可呈示范畴。

\section{Betti 局部系统与 de Rham 局部系统}\label{sec:11.2}\index{局部系统!Betti 与 de Rham 比较|hyperpage}\index{Riemann–Hilbert 对应|hyperpage}\index{刚性化|hyperpage}
在 de Rham 情形，谱参数是 de Rham 局部系统：
\[\begin{split}
\LocSys^{\mathrm{dR}}_{\checkG}(X)&=\operatorname{Conn}_{\checkG}(X)\\
&:=\{(E_{\checkG},\nabla)\mid E_{\checkG}\text{ 为 }\checkG\text{-挠子，}\nabla\text{ 为联络}\}.
\end{split}\]
在 Betti 情形，谱参数是 Betti 局部系统：
\[\LocSys^{\mathrm{Betti}}_{\checkG}(X):=\operatorname{Map}(X^{\mathrm{top}},B\checkG).\]
选定基点 $x\in X(\mathbb C)$ 后，刚性化 Betti 模空间的经典截断为
\[\LocSys^{\mathrm{Betti},\mathrm{rig},\mathrm{cl}}_{\checkG,x}(X)\simeq\Hom(\pi_1(X(\mathbb C),x),\checkG).\]
未刚性化的经典叠通过取共轭作用的商得到。这个描述只是第一近似：范畴 Langlands 中使用的 Betti 模空间是导出的，它的切复形由 $X(\mathbb C)$ 上以伴随局部系统为系数的上链控制。

这两个模叠在代数意义上并不相同。de Rham 叠依赖于复代数曲线 $X$，而 Betti 叠只依赖于 $X(\mathbb C)$ 的拓扑。
\begin{example}
令 $X=E$ 为椭圆曲线，并选取 $e\in E$。则
\[\LocSys^{\mathrm{dR},\mathrm{rig}}_{\Gm,e}(E)\longrightarrow\Pic_E^0\simeq E^\vee\]
是普遍向量扩张，因此是 $E^\vee$ 上的 $\Ga$-挠子，等价地，是一个 $\mathbb A^1$-丛。相比之下，
\[\LocSys^{\mathrm{Betti},\mathrm{rig}}_{\Gm,e}(E)\simeq\Hom(\mathbb Z^2,\Gm)\simeq\Gm^2.\]
\srcpage{58}
Riemann–Hilbert 对应在解析意义下识别它们的复点，但这一对应使用指数映射，并不识别代数结构。
\end{example}

\section{为何朴素的 Betti 猜想过大}\label{sec:11.3}\index{拟光滑层|hyperpage}\index{Fourier–Mellin 变换|hyperpage}
最初可能猜想
\[\Shv^{\mathrm{all}}(\Bun_G)\overset{?}{\simeq}\IndCoh_{\mathrm{Nilp}}(\LocSys^{\mathrm{Betti}}_{\checkG}).\]
\begin{example}
即使 $G=\GL_1$、$X=E$ 为椭圆曲线，这个猜想也是错误的。在自守侧，
\[\Bun_{\Gm}\simeq\coprod_{d\in\mathbb Z}\Pic_E^d\times B\Gm\simeq E\times\mathbb Z\times B\Gm.\]
拓扑 Fourier–Mellin 变换将有限性或可构造性的归一化识别为
\[\Shv(\mathbb Z)_{\mathrm{fin}}\simeq\QCoh(B\Gm),\qquad
\Shv(B\Gm)_{\mathrm{fin}}\simeq\QCoh(\mathrm{pt}\times_{\mathbb A^1}\mathrm{pt}).\]
若按字面取所有层，则这些范畴需要完备化；在这个例子中，完备化问题并不是两侧不匹配的原因。秩一局部系统的导出 Betti 叠具有相应的因子：
\[\LocSys^{\mathrm{Betti}}_{\Gm}(E)\simeq\Gm^2\times B\Gm\times(\mathrm{pt}\times_{\mathbb A^1}\mathrm{pt}).\]
因此，唯一不匹配的是椭圆曲线因子：
\[\Shv^{\mathrm{all}}(E)\ne\QCoh(\Gm^2).\]
正确的替代是光滑层组成的满子范畴
\[\begin{split}
\operatorname{QLisse}(E):=\{\cF\in\Shv^{\mathrm{all}}(E)\mid{}
& H^i(\cF)\text{ 是由有限维局部系统}\\[-2pt]
&\text{生成的 ind-对象}\}.
\end{split}\]
等价地，
\[\operatorname{QLisse}(E)=\Shv^{\mathrm{all}}_0(E):=\{\cF\mid\operatorname{SS}(\cF)\subset T_E^*E\}.\]
这里 $T_E^*E$ 是零截面。因此，朴素猜想的错误恰好在于包含了具有非零微局部支撑的层。
\end{example}

\section{层的微局部奇异支撑}\label{sec:11.4}\index{奇异支撑!微局部|hyperpage}\index{特征簇|hyperpage}\index{拟光滑层|hyperpage}
令 $S$ 为光滑复概形。对可构造层 $\cF$，微局部分析赋予一个闭锥形子集\cite{Bei16}
\[\operatorname{SS}(\cF)\subset T^*S,\]
称为奇异支撑。它度量 $\cF$ 不能局部常值的程度。更准确地说，若切方向 $v\in T_sS$ 被 $\operatorname{SS}(\cF)_s$ 中的所有余向量消去，则 $\cF$ 沿这一方向局部无圈，或者说是光滑的。对闭锥形子集 $N\subset T^*S$，通过 ind-完备化定义所有层的版本：
\begin{definition*}{}
\[\Shv^{\mathrm{all}}_N(S):=\operatorname{Ind}\{\cF\in\Shv_{\mathrm{cons}}(S)\mid\operatorname{SS}(\cF)\subset N\}.\]
\end{definition*}
等价地，也可以将它看作 $\Shv^{\mathrm{all}}(S)$ 中由奇异支撑位于 $N$ 的可构造对象生成的余完备满子范畴。零支撑范畴就是光滑层范畴，也称拟光滑层范畴：
\[\Shv^{\mathrm{all}}_0(S)\simeq\operatorname{QLisse}(S).\]
在 AGKRRV 的归一化中，这指的是有限秩局部系统的 ind-范畴的左完备化。

\begin{transnote}{易错预警}
「左完备化」是相对 $t$-结构取的（Lurie SAG 意义），不是 $\operatorname{Ind}$-完备化：它只把沿 $t$-结构正向收敛的对象补进来，补完恰是零奇异支撑层的完备化，即拟光滑（quasi-lisse）层。若按普通 $\operatorname{Ind}$ 去补，会把不满足收敛条件的对象一并放进来，微局部意义就丢了。[来源：nLab "t-structure"，https://ncatlab.org/nlab/show/t-structure（Last revised 2023-04-20；2026-11-17 访问）；并见原书 p.58 本段。]
\end{transnote}

为了与 de Rham 侧比较，若 $M$ 为正则完整 $D$ 模，而 $\operatorname{Sol}(M)$ 是它的解层，则
\[\operatorname{SS}(\operatorname{Sol}(M))=\operatorname{Char}(M).\]
\srcpage{59}
因此，Betti 层的微局部支撑是完整 $D$ 模的特征簇的 Betti 对应物。

\section{全局幂零锥}\label{sec:11.5}\index{全局幂零锥|hyperpage}\index{Hitchin 纤维化|hyperpage}
对 $G$-丛 $P$，记 $\mathfrak g_P=P\times^G\mathfrak g$。$\Bun_G$ 的经典余切叠是 Higgs 叠
\[T^*\Bun_G\simeq\operatorname{Higgs}_G,\qquad
\operatorname{Higgs}_G(\mathbb C)=\{(P,\varphi)\mid\varphi\in H^0(X,\mathfrak g_P^*\otimes\Omega_X)\}.\]
\begin{definition*}{}
\[\mathrm{Nilp}_{\mathrm{glob}}:=\{(P,\varphi)\in T^*\Bun_G\mid
\varphi_x\in\mathcal N_{\mathfrak g^*}\text{ 对每个 }x\in X\text{ 都成立}\}.\]
\end{definition*}
它是通常的幂零锥 $\mathcal N_{\mathfrak g^*}\subset\mathfrak g^*$ 的全局对应物。

等价地，$\mathrm{Nilp}_{\mathrm{glob}}$ 是 Hitchin 纤维化的零纤维。若 $p_i\in\mathbb C[\mathfrak g^*]^G$ 是次数为 $d_i$ 的齐次生成元，置
\[\cA_G:=\bigoplus_i H^0(X,\Omega_X^{d_i}).\]
Hitchin 映射为
\[h:T^*\Bun_G\longrightarrow\cA_G,\qquad(P,\varphi)\longmapsto(p_i(\varphi))_i.\]
于是
\[\mathrm{Nilp}_{\mathrm{glob}}=h^{-1}(0).\]
用不变记号表示，这就是映射
\[\Gamma(X,\mathfrak g_P^*\otimes\Omega_X)\longrightarrow\Gamma(X,(\mathfrak g^*\!/\!/ G)_{\Omega_X}),\]

\begin{transnote}{含糊点拨}
扭转 $(\text{--})_{\Omega_X}$ 把权重为 $d_i$ 的不变量与 $\Omega_X^{\otimes d_i}$ 张量：$\mathbb C[\mathfrak g^*]^{\check G}$ 的齐次生成元 $p_i$ 次数为 $d_i$，所以不变量截面是 $\bigoplus_i H^0(X,\Omega_X^{d_i})$，而不是朴素的 $H^0(X,(\mathfrak g^*\!/\!/\check G)\otimes\Omega_X)$。必须逐生成元各自扭转，因为各生成元权不同，不能先取商再统一扭。这个直和正是 Hitchin 基 $\cA_{\check G}$。[依据：原书 p.59 本段；并见\chref{14}开头对 Hitchin 底空间的直和描述。]
\end{transnote}
其中 $(\text{--})_{\Omega_X}$ 表示由 $\mathfrak g^*$ 上的伸缩作用产生的 $\Gm$-扭转。取次数为 $d_i$ 的齐次生成元 $p_i$ 后，该目标为 $\bigoplus_iH^0(X,\Omega_X^{d_i})$。

\section{具有幂零支撑的 Betti 范畴猜想}\label{sec:11.6}\index{Betti 几何 Langlands|hyperpage}\index{紧生成|hyperpage}\index{Hecke 光滑性|hyperpage}
修正后的 Betti 猜想是 Betti 几何 Langlands 的幂零支撑版本\cite{BZN16,NY19}：
\begin{conjecture*}{}
\[\Shv^{\mathrm{all}}_{\mathrm{Nilp}}(\Bun_G)\simeq\IndCoh_{\mathrm{Nilp}}(\LocSys^{\mathrm{Betti}}_{\checkG}),\]
其中
\[\Shv^{\mathrm{all}}_{\mathrm{Nilp}}(\Bun_G):=\{\cF\in\Shv^{\mathrm{all}}(\Bun_G)\mid\operatorname{SS}(\cF)\subset\mathrm{Nilp}_{\mathrm{glob}}\}.\]
\end{conjecture*}
在适用于一般层理论，特别是 $\ell$-进层的表述中，谱叠被替换为受限变体 $\LocSys^{\mathrm{restr}}_{\checkG}$（定义见 \secref{12.5}），但自守范畴仍然是 $\Bun_G$ 上层范畴的幂零支撑子范畴。

这个条件通过三项基本检验。第一，它对环面成立：由于 $\mathfrak t^*$ 中的幂零锥是零截面，$\Shv^{\mathrm{all}}_{\mathrm{Nilp}}(\Bun_T)$ 就是光滑部分，而 Mellin 变换将它识别为 Betti 对偶环面局部系统叠上的拟凝聚层。第二，它不会过小：Hecke 特征层应当具有幂零奇异支撑。事实上，Nadler–Yun 证明，对具有幂零奇异支撑的层，Hecke 函子在 $X$ 的点这一参数上局部常值\cite{NY19}；结合 AGKRRV 中使用的逆命题和推广\cite{AGK20a}，Hecke 光滑行为等价于幂零奇异支撑。特别地，特征层属于 $\Shv^{\mathrm{all}}_{\mathrm{Nilp}}(\Bun_G)$。第三，它不会过大：$\Shv^{\mathrm{all}}_{\mathrm{Nilp}}(\Bun_G)$ 是紧生成的，而 $\Shv^{\mathrm{all}}(\Bun_G)$ 一般不是。

更一般地，对光滑概形 $S$ 和锥形子集 $N\subset T^*S$，范畴 $\Shv^{\mathrm{all}}_N(S)$ 受一个适配于 $N$ 的分层控制：奇异支撑位于 $N$ 的对象沿各个分层局部常值。这就是幂零支撑恢复有限性性质的微局部原因。
